\documentclass[10pt,a4paper,roman,english,colorlinks]{moderncv}
\moderncvstyle{classic}
\moderncvcolor{blue}

\usepackage[utf8]{inputenc}
\usepackage[scale=0.86,a4paper]{geometry}
\usepackage[english]{babel}
\setlength{\hintscolumnwidth}{3cm}

\firstname{Anjali}
\familyname{Suman}
\title{Data Science \& Statistics}
\email{anjali.suman.stats@gmail.com}
\homepage{github.com/AnjaliML}

\begin{document}
\hypersetup{urlcolor=blue}
\makecvtitle

\section{Profile}
\cvitem{}{Statistics graduate with experience in applied analysis, Python projects and research support. Currently studying data science with machine learning and AI through Cambridge PACE and FourthRev. Work includes customer segmentation, regression implementation, Bayesian optimisation and reproducible learning resources.}

\section{Technical Skills}
\cvitem{Programming}{Python, R, MATLAB; working knowledge of C/C++}
\cvitem{Data \& ML}{pandas, NumPy, scikit-learn, PyTorch, statistical modelling, data visualisation}
\cvitem{Tools}{Jupyter, Git/GitHub, \LaTeX}

\section{Selected Projects}
\cventry{Oct 2026}{E-commerce Customer Segmentation}{Cambridge Career Accelerator mini-project}{}{}{Analysed 951,669 transaction rows in Python and engineered five features for 68,300 customers. Checked data quality, fitted five K-means segments and compared cluster counts using inertia and sampled silhouette analysis. Profiled segments in original units and visualised them with PCA and t-SNE; used a sampled hierarchical dendrogram as exploratory evidence.}
\cventry{2026}{\href{https://github.com/AnjaliML/DataScience-Python101}{Data Science Python 101}}{Python learning resource}{}{}{Lessons, exercises, automated checks and accessibility features. Covers data preparation, analysis and model evaluation.}
\cventry{2024--25}{\href{https://github.com/AnjaliML/ML101}{ML101}}{Machine learning notebooks}{}{}{Implemented examples of Gaussian process regression, Bayesian optimisation and related methods.}
\cventry{2024}{\href{https://github.com/AnjaliML/findingPi}{findingPi}}{Numerical computing notebooks}{}{}{Explored algorithms for estimating $\pi$ and plotted their convergence.}
\cvitem{Earlier work}{\href{https://github.com/AnjaliML/InterStateMigrationIndia}{Inter-state migration analysis} (2022); \href{https://github.com/AnjaliML/FantasyCricket}{fantasy-cricket score modelling} (2017).}

\section{Research \& Experience}
\cventry{Acknowledged 2025}{Regression implementation support}{Research contribution}{}{}{Helped implement a Python regression algorithm to estimate model prefactors for drop-impact research. Named in the acknowledgements of V. Sanjay and D. Lohse, \href{https://doi.org/10.1103/PhysRevLett.134.104003}{Physical Review Letters 134, 104003 (2025)}.}
\cventry{2022--Apr 2025}{Research collaboration}{Oklahoma State University (remote)}{}{}{Investigated temporal Bayesian optimisation and machine-failure prediction using survival analysis.}
\cventry{2019--20}{Design \& Analysis}{Sureland Industrial Fire Safety Limited}{}{}{Analysed fire-incident data and smoke-detector placement.}
\cventry{Oct 2026--present}{Learner representative}{Cambridge PACE / FourthRev programme}{}{}{Participated with fellow representatives in the 21 October course-team meeting, raising questions about Python support and assessment guidance and contributing learners' perspectives.}

\section{Education \& Development}
\cventry{2026--present}{Data Science with Machine Learning \& AI Career Accelerator}{University of Cambridge Professional and Continuing Education / FourthRev}{}{}{In progress; began August 2026. Applied statistics, Python and data-science projects.}
\cvitem{2026}{\href{https://www.futurelearn.com/certificates/svp37fa}{Understanding the GDPR}, University of Groningen via FutureLearn; Certificate of Achievement, 97\%.}
\cventry{2018}{M.Sc. Statistics}{Hindu College, University of Delhi}{}{}{}
\cventry{2016}{B.Sc. (Hons) Statistics}{Ramanujan College, University of Delhi}{}{}{}

\end{document}
